\section{Coleta de lixo}
\label{sec:gc}

Em Java você aloca e o \termo{garbage collector}{Coletor de lixo (GC):
componente da JVM que descobre quais objetos não são mais alcançáveis e
recupera a memória deles.} recupera. A escolha do coletor é uma troca
entre três coisas: \emph{throughput} (fração do tempo gasta no seu
código), latência (duração das pausas) e memória. O guia de GC da Oracle
para o JDK~25 (\citetitle{doc:gc-collectors}) descreve os coletores da
HotSpot assim:

\small
\begin{longtable}{@{}p{1.8cm} p{3.2cm} p{7.2cm}@{}}
\toprule
Coletor & Opção & Segundo a documentação \\
\midrule \endhead
Serial & \texttt{-XX:+UseSerialGC} & Uma thread faz todo o trabalho de
  GC; indicado para dados pequenos (até cerca de 100~MB) ou uma CPU \\
Parallel & \texttt{-XX:+UseParallelGC} & Como o Serial, mas com várias
  threads; para quando throughput máximo importa e pausas de um segundo
  ou mais são aceitáveis \\
G1 & \texttt{-XX:+UseG1GC} & Majoritariamente concorrente; busca cumprir
  uma meta de pausa com alta probabilidade mantendo bom throughput. É o
  padrão (\citetitle{doc:jep248}, JDK~9) \\
ZGC & \texttt{-XX:+UseZGC} & Baixa latência, pausas independentes do
  tamanho do heap; o guia fala em pausas máximas abaixo de um
  milissegundo, ao custo de algum throughput \\
\bottomrule
\end{longtable}
\normalsize

\paragraph{A evolução do ZGC.}
O ZGC ficou pronto para produção no \citetitle{doc:jep377} (JDK~15).
O \citetitle{doc:jep439} (JDK~21) acrescentou o modo geracional, o
\citetitle{doc:jep474} (JDK~23) o tornou padrão e o
\citetitle{doc:jep490} (JDK~24) removeu o modo não geracional. Em JDK~25,
\texttt{-XX:+UseZGC} é sempre o ZGC geracional; quem ainda passa
\texttt{ZGenerational} recebe aviso de opção obsoleta.

\paragraph{Shenandoah.}
O \citetitle{doc:jep379} (JDK~15) tornou o Shenandoah um coletor de
produção no OpenJDK; a página do projeto (\citetitle{doc:shenandoah}) o
descreve como um coletor de pausas ultrabaixas que faz mais trabalho
concorrente com o programa. Ele não aparece no guia de GC da Oracle
citado acima: confirme se a distribuição de JDK que você usa o inclui
antes de planejar com ele.

\paragraph{A decisão que isso muda.}
O guia da Oracle recomenda começar sem escolher nada: deixe a JVM
escolher, ajuste o tamanho do heap e só então troque de coletor, se o
requisito não for atendido. Na prática:
\begin{itemize}
\item API com requisito de latência (p99) apertado: compare G1 e ZGC
  medindo o p99 no teste de carga da Parte~I, não a média.
\item Job batch que só precisa terminar rápido: Parallel é candidato.
\item Container pequeno, com pouca memória e uma CPU: Serial é o que o
  guia indica para esse perfil.
\end{itemize}
Antes de trocar de GC, confira se o problema não é alocação demais: o
autoboxing da seção~\ref{sec:autoboxing} é um jeito comum de criar
objetos sem perceber. O JFR (seção~\ref{sec:ferramentas}) mostra pausas
e alocação por método.
